%! AP Statistics — Unit 7, Lesson 7.2 — Guided Notes
\documentclass[10pt]{article}
\usepackage{apstats-article}
\usepackage{apstats-boxes}
\newcommand{\TallMath}[1]{$\displaystyle #1\rule[-1.4em]{0pt}{3.2em}$}

\begin{document}

\pageheader{Unit 7, Lesson 7.2}{Guided Notes}
\namedateperiod

\begin{objectivebox}
By the end of this lesson, I will be able to:
\begin{itemize}
    \item[$\square$] \writeline
    \item[$\square$] \writeline
    \item[$\square$] \writeline
    \item[$\square$] \writeline
\end{itemize}
\end{objectivebox}

\begin{vocabbox}
\termblanklong{$t$-distribution}
\termblanklong{Degrees of freedom ($df$)}
\termblanklong{Standard error of $\bar x$ ($SE_{\bar x}$)}
\termblanklong{Critical value $t^*$}
\termblanklong{One-sample $t$-interval}
\termblanklong{Matched-pairs design}
\end{vocabbox}

\begin{hookbox}
A nutritionist records daily sodium intake for a random sample of 20 college students:
$\bar x = 2840$ mg, $s = 480$ mg. The recommended maximum is 2300 mg.

Why can't we use a $z$-interval to estimate the true mean sodium intake? \writelines{2}
\end{hookbox}

\begin{notesbox}{The $t$-Distribution (VAR-7.A)}
When we substitute $s$ for $\sigma$, the statistic $t = \dfrac{\bar x - \mu}{s/\sqrt{n}}$
follows a \blank{3cm}-distribution.

\vspace{0.1in}
\textbf{Properties:}
\begin{itemize}
    \item Bell-shaped, \blank{4cm}, centered at \blank{2cm}.
    \item Heavier tails than the standard normal (because $s$ adds variability).
    \item Indexed by \textbf{degrees of freedom}: $df = n - $ \blank{1cm}.
    \item As $df$ \blank{4cm}, the $t$-distribution approaches $N(0, 1)$.
\end{itemize}

\textbf{Finding $t^*$:} For 95\% CI with $n = 20$: $df = $ \blank{2cm}, $t^* = $ \blank{2.5cm}.
\end{notesbox}

\begin{notesbox}{Conditions for a One-Sample $t$-Interval (UNC-4.P)}
\begin{enumerate}
    \item \textbf{Random:} Data come from a \blank{5cm} or randomized experiment.
    \item \textbf{10\%:} If sampling w/o replacement: $n \le $ \blank{2cm}$\% \cdot N$.
    \item \textbf{Large Sample:}
          \begin{itemize}
              \item $n \ge $ \blank{2cm}: CLT applies $\to$ condition met.
              \item $n < $ \blank{2cm}: check that the data show no strong
                    \blank{4cm} or \blank{3cm}.
          \end{itemize}
\end{enumerate}
\end{notesbox}

\begin{notesbox}{The One-Sample $t$-Interval Formula (UNC-4.Q, UNC-4.R)}
\renewcommand{\arraystretch}{2}
\begin{tabularx}{\linewidth}{lX}
    Point estimate & $\bar x$ \\
    Standard error & $SE_{\bar x} = \dfrac{s}{\sqrt{n}}$ \\
    Margin of error & $ME = t^* \cdot \dfrac{s}{\sqrt{n}}$ \quad ($t^*$ from table, $df = n - 1$) \\
    \textbf{Interval} & $\bar x \pm t^* \cdot \dfrac{s}{\sqrt{n}}$ \\
\end{tabularx}

\vspace{0.1in}
\textbf{Sodium example:} $n = 20$, $\bar x = 2840$, $s = 480$, 95\% CI.

$df = $ \blank{2cm}, \quad $t^* = $ \blank{2.5cm}, \quad $SE = 480/\sqrt{20} = $ \blank{2.5cm} mg

Interval: $2840 \pm $ \blank{2cm} $\cdot$ \blank{2cm} $= 2840 \pm$ \blank{2cm}
$= ($ \blank{2cm}$,$ \blank{2cm}$)$ mg
\end{notesbox}

\begin{notesbox}{The Four-Step Procedure (PCCF)}
\begin{enumerate}
    \item \textbf{Parameter:} Let $\mu = $ the true mean \blank{4cm} for \blank{4cm}.
    \item \textbf{Conditions:} State and verify all three (Random, 10\%, Large Sample).
    \item \textbf{Calculate:} $\bar x \pm t^* \cdot s/\sqrt{n}$
    \item \textbf{Conclude:} ``We are \blank{2cm}\% confident that the interval
          $(\rule{1.5cm}{0.4pt},\ \rule{1.5cm}{0.4pt})$ captures the true mean
          \blank{4cm} for \blank{4cm}.''
\end{enumerate}
\end{notesbox}

\begin{notesbox}{Matched-Pairs Design (UNC-4.O.3)}
When data come in pairs (before/after, two treatments on same subject):
\begin{enumerate}
    \item Compute differences $d_i = x_{i,1} - x_{i,2}$ for each pair.
          (State the order of subtraction: \writeline)
    \item Find $\bar d = $ \blank{3cm} and $s_d = $ \blank{3cm}.
    \item Apply the one-sample $t$-interval to the differences:
          $\bar d \pm t^* \cdot \dfrac{s_d}{\sqrt{n}}$, \quad $df = n - 1$,
          where $n = $ number of \blank{3cm}.
\end{enumerate}
\end{notesbox}

\end{document}
